%% 30-methodology.tex — Methodology

\section{Methodology}
\label{sec:methodology}

\subsection{Corpus-first Principle}

The alignment methodology follows a single governing principle: \emph{profile
the corpus before constructing any mapping rule}.
Schema-level alignment presupposes that the source ontology specification
faithfully describes the semantics of deployed data.
In the DDB-EDM case it does not: the specification defines \texttt{dc:type}
as a free-text or controlled-vocabulary property, but says nothing about
which controlled vocabulary is used in each sector, or how consistently it
is applied.
The corpus-first step produces a field profile --- a frequency table of every
(entity type, JSON key) pair observed in the data --- which reveals which
fields actually occur, at what coverage, and with what value distributions.
Alignment rules are constructed against this profile, not against the
specification.

The pilot corpus used in this study consists of 115{,}432 DDB records
returned by the DDB API for the search terms \textit{Goethe} and
\textit{Faust} \todo{confirm exact query and date of retrieval}.
Despite its thematic focus, the corpus spans all seven DDB institutional
sectors (Archiv, Bibliothek, Denkmal, Forschung, Mediathek, Museum, Sonstige)
and all five mediatypes (Audio, Foto, Text, Video, Nicht digitalisiert),
making it representative of the structural diversity of DDB-EDM data.
Each record is a JSON-LD document as returned by the DDB Search
API~\todo{cite DDB API documentation}, structured around the EDM entities
\texttt{edm:ProvidedCHO}, \texttt{edm:WebResource}, \texttt{edm:Aggregation},
and \texttt{edm:PhysicalThing}.

\subsection{Overview: Three-Layer Dispatch}

The transformation pipeline
\texttt{transform\_edm\_to\_mocho.py} applies three independent layers to
each record, combining their outputs into a single N-Triples stream
(Figure~\ref{fig:pipeline}).

\begin{figure}[t]
\centering
\todo{Pipeline figure: four-box flowchart
  [JSON-LD record] → [Layer 1: Property alignment]
  → [Layer 2: htype dispatch] → [Layer 3: dc:type dispatch] → [N-Triples].
  Suitable for a single-column figure at 0.9\textbackslash textwidth.}
\caption{Three-layer dispatch pipeline. Layers 2 and 3 are independent and
  can both fire for the same record.}
\label{fig:pipeline}
\end{figure}

\noindent
The three layers address orthogonal sub-problems:
Layer~1 maps \emph{property} IRIs (DC $\to$ RDA);
Layer~2 assigns \emph{archival structural type} from a DDB-specific hierarchy
code;
Layer~3 assigns \emph{domain WEMI class} from the free-text or
controlled-vocabulary document type.
Layers~2 and~3 are independent: both can fire for the same record, and their
outputs are unioned in the triple store.

\subsection{Layer 1: Property Alignment}

\noindent
\textbf{Input.} For each EDM entity in a record (ProvidedCHO, WebResource,
Aggregation, Agent, Place, TimeSpan), every JSON key is looked up in the
alignment table \texttt{alignment\_ddbedm\_mocho.csv}.

\noindent
\textbf{Alignment table construction.}
The table was built in two steps:
(1) a corpus profiling script enumerates all (entity type, JSON key) pairs
observed across the 115{,}432 records, recording occurrence frequency and
percentage coverage;
(2) each JSON key is resolved to its EDM/DC IRI, then matched to an RDA
sub-property IRI via the \texttt{mapping\_dct\_to\_rda.csv} file~\todo{cite
or describe provenance of this mapping --- derived from RDA Registry
\texttt{mapElement2DCT.xml}?}.
Only rows where the \texttt{in\_mocho} flag is \texttt{True} are used;
rows without an RDA sub-property in the mocho OWL import are documented as
known gaps (Section~\ref{subsec:gaps}).

\noindent
\textbf{Special-case keys.}
Three high-fanout properties bypass the table and receive fixed predicates:
\texttt{dc:creator} maps to the generic \texttt{rdam:P30263} (has creator
agent of manifestation, 232 RDA sub-properties);
\texttt{dc:contributor} maps to \texttt{dc:contributor} (no generic RDA
equivalent in the current mocho import);
and \texttt{dcSubject}/\texttt{dcTermsSubject}/\texttt{dcTermSubject} are
deduplicated before emit, routing URI values to \texttt{dcterms:subject} and
literals to \texttt{dc:subject}.

\noindent
\textbf{EDM-structural properties.}
Properties carrying provenance and access metadata
(\texttt{edm:isShownAt}, \texttt{edm:aggregatedCHO},
\texttt{edm:dataProvider}, \texttt{edm:hasView}, \texttt{edm:isShownBy},
\texttt{edm:provider}) are explicitly classified as out of scope for
DC\,$\to$\,RDA routing and ignored in Layer~1.
These properties require a separate EDM-to-mocho direct mapping pass and are
not emitted in the current pipeline; their exclusion is documented in the
alignment table rather than silently dropped.

\subsection{Layer 2: Archival Hierarchy Typing}

DDB archival records carry a \texttt{ddb:hierarchyType} code on both
\texttt{edm:PhysicalThing} (archival containers) and \texttt{edm:ProvidedCHO}
entities, encoding nine structural levels of an archive finding aid:
from repository (\textit{Tektonische Sammlung}) down to individual archival
document (\textit{Archivale}) and its parts (\textit{Teil}).

Layer~2 looks up each observed \texttt{hierarchyType} code in
\texttt{lookup\_htype\_doco\_rico.csv} and emits two triples:
an \texttt{rdf:type} assertion using the corresponding RiC-O class
(\texttt{rico:RecordSet}, \texttt{rico:Record}, or \texttt{rico:RecordPart}),
and a \texttt{rico:hasRecordSetType} triple pointing to both the standard
RiC-RST individual (e.g.\ \texttt{ric-rst:Fonds}) and a DDB-specific
vocabulary term (\texttt{vocnet-htype:ht030}) for the finer-grained level.
Containment relations between records use \texttt{rico:includesOrIncluded};
a \texttt{rico:Record} containing a \texttt{rico:RecordPart} uses
\texttt{rico:hasOrHadConstituent}.

\texttt{PhysicalThing} ancestors receive the RiC-O class only; no
\texttt{mocho:Manifestation} is emitted for container nodes.
The \texttt{ProvidedCHO} entity receives both the RiC-O type (from Layer~2)
and the \texttt{mocho:Manifestation} plus dc:type dispatch (from Layer~3),
enabling cross-layer queries over the same node.

\subsection{Layer 3: Document-Type Dispatch}

\noindent
\textbf{Lookup key.}
The dispatch table \texttt{lookup\_dctype\_to\_class.csv} is keyed on
(\textit{mediatype}, \textit{sector}, \textit{dc\_type\_de}), where mediatype
and sector are read from the record's \texttt{Concept[]} array
(or from \texttt{ProvidedCHO.edmType} as a simpler mediatype path),
and \texttt{dc\_type\_de} is the German free-text or controlled-vocabulary
value of \texttt{dc:type} (e.g.\ \textit{Fotografie}, \textit{Zeichnung}).

\noindent
\textbf{Three-level fallback.}
The lookup proceeds through three levels:
(i) exact match on (mediatype, sector, dc\_type\_de);
(ii) any-sector match on (mediatype, any, dc\_type\_de);
(iii) any-mediatype match on (any, any, dc\_type\_de).
If no match is found at any level, the fallback \texttt{mocho:Manifestation}
is asserted.

\noindent
\textbf{WEMI-slot semantics.}
The dispatch table has two class slots per row (W and M):
\begin{itemize}
  \item A \textbf{W-slot} match (Work-level class, e.g.\ \texttt{vra:Work}
        for a museum photograph) is emitted \emph{instead of}
        \texttt{mocho:Manifestation} for the ProvidedCHO.
  \item An \textbf{M-slot} match (Manifestation-level class) is emitted
        \emph{alongside} \texttt{mocho:Manifestation}.
  \item No match defaults to \texttt{mocho:Manifestation} only.
\end{itemize}

This mirrors \textsc{mocho}'s \texttt{owl:unionOf} design: the WEMI-level
class is valid only if it belongs to at least one domain-specific
class~\cite{tan2026mocho}.

\noindent
\textbf{WebResource typing (mt002 / Photo).}
For photo-type records, each \texttt{edm:WebResource} URI additionally
receives:
\texttt{mocho:ImageObject}, \texttt{rdac:C10007} (RDA Manifestation),
\texttt{rdam:P30001 rdact:1018} (carrier type: online resource), and
\texttt{vra:imageOf} pointing back to the ProvidedCHO URI.
This makes the digital surrogate independently queryable as a manifestation
of the physical object.
